\input{../../../stock/en-tete_v6.tex}

\newcommand{\titrechapitre}{Cours d'introduction}
\newcommand{\numerochapitre}{1}

\renewcommand{\headrulewidth}{0.4pt}
\renewcommand{\footrulewidth}{0.4pt}
\fancyfoot[L]{Notes de Raphaël JONTEF}
\fancyfoot[C]{\thepage}
\fancyfoot[R]{Enseirb-Matmeca}
\fancyhead[L]{Cours de Logique et preuve}
\fancyhead[R]{\titrechapitre}

\begin{document}
\input{../../../stock/commands.tex}

% Page de titre custom
\setlength{\headheight}{15pt}
\begin{titlepage}
    \centering
    \vspace*{3cm}
    {\huge Chapitre \numerochapitre\par}
    {\Huge \bfseries\titrechapitre\par}
    \vspace{2cm}
    {\Large Notes rédigées par Raphaël JONTEF\par}
    \vspace{0.5cm}
    {\normalsize Cours enseigné par Frédéric HERBRETEAU\par}
    \vspace{4cm}
    {\small Enseirb-Matmeca\par}
    {\small filière informatique\par}
    \vfill
    {\large \today\par}
\end{titlepage}

% plan du cours
\tableofcontents
\newpage


% -----------Corps du document--------------

\section*{Sommaire des chapitres}

Les chapitres se présenteront dans l'ordre suivant~:
\begin{enumeratebf}
    \item* Technique de preuve (rappels en TD)
    \item Introduction
    \item Logique propositionnelle
    \item Introduction à la logique des prédicats
    \item Programmation logique, IA, solveurs
    \item Induction structurelle
    \item Ensemble bien fondés et terminaison
    \item Machines à états et invariants
    \item Machines à états et invariants
    \item Spécification et preuve de programmes
\end{enumeratebf}


8h de cours de 1h20 chacuns, 7 séances de TD de 2h chacunes

cours ici : https://thor.enseirb-matmeca.fr/sites/if107-main/

\section{Introduction} 

\begin{definition}{}{proposition}
    Une \notion{proposition} est un énoncé qui est soit vrai, soit faux.
\end{definition}

\begin{exemple}{}{}
    \begin{itemize}
        \item $2+3 = 5$ est une proposition vraie.
        \item "Pour tout $n \in \N$", $n^2 + n + 41$ est un nombre premier" est fausse (prendre $n=40$).
    \end{itemize}
\end{exemple}

\begin{definition}{}{axiome}
    Un \notion{axiome} est une proposition supposée vraie.
\end{definition}

\begin{exemple}{}{}
    \begin{itemize}
        \item La transitivité de la relation $\leq$ sur la classe des entiers relatifs.
        \item Par toute paire de points passe une unique droite.
    \end{itemize}
\end{exemple}

\begin{definition}{}{théorie}
    Un ensemble d'axiomes définit une \notion{théorie}, qui fixe un cadre de raisonnement.
\end{definition}

\begin{exemple}{}{}
    La théorie Euclidienne est une théorie de la géométrie basée sur les axiomes d'Euclide.
\end{exemple}

\begin{exemple}{}{}
    En géométrie ellpitique~:
    \begin{itemize}
        \item étant donnée une droite $(D)$ et un point $A$ n'appartenant pas à $(D)$, il n'existe aucune droite parallaèle à $(D)$ passant par $A$.
        \item la somme des anles d'un triangles faire entre 180 et 540 degrés.
    \end{itemize}
\end{exemple}

\begin{definition}{}{consistance d'une théorie}
    Une théorie est \notion{consistante} si elle ne permet pas de démontrer une propriété et son contraire (contradiction).
\end{definition}

\begin{definition}{}{complétude d'une théorie}
    Une théorie est \notion{complète} si pour toute proposition $P$, soit $P$ est démontrable, soit sa négation $\neg P$ est démontrable.
\end{definition}


\begin{remarque}{}{}
    une théorie peut être consistante sans être complète (exemple : arithmétique de Peano)
\end{remarque}

\begin{definition}{}{preuve}
    On appelle \notion{preuve} la succession d'étapes de déductions, partant d'axiomes et aboutissant à la démonstration d'une proposition.
\end{definition}





\end{document}
